\section{Model training and experiment tracking}
\label{sec:mlflow}

Stage 7 does all the training. Table~\ref{tab:training} records the settings, all read from
\texttt{config.yaml} rather than typed into the notebook.

\begin{table}[htbp]
\centering
\caption{Training configuration.}
\label{tab:training}
\small
\begin{tabular}{|p{4.0cm}|p{10.0cm}|}
\hline
\textbf{Setting} & \textbf{Value} \\
\hline
Modelling frame & 545 games, at least 200 English early-access reviews and 50 post-launch \\
\hline
Held-out cohort & 424 games still in early access, never used for training \\
\hline
Split & temporal on launch date: 381 train, 164 test \\
\hline
Random seed & 42 \\
\hline
Predictors & 19, an explicit list rather than every column minus the outcome \\
\hline
Significance level & 0.01667, Bonferroni-corrected over three hypotheses \\
\hline
Band cutoffs & 43.6 and 57.2, from training terciles, written back to \texttt{config.yaml} \\
\hline
\end{tabular}
\end{table}

Every fitted run is logged to MLflow. Version 3 dropped the plain file store, so the record is a
SQLite database rather than a folder of files, which changes only the command used to open it:

\begin{cmd}
mlflow ui --backend-store-uri sqlite:///data/ml_runs/mlflow.db
\end{cmd}

\noindent Then open \url{http://127.0.0.1:5000}; Ctrl+C stops it. Note the forward slashes. The
\texttt{sqlite:///} prefix is a URI, not a Windows path, and backslashes are rejected there even on
Windows.

\shot{mlflow_runs}{0.98}{The run list. Each row is one fitted model with the metrics it scored.}

\shot{mlflow_run_detail}{0.72}{One run in detail. The parameters record the frame, the split and
the seed; the metrics record what it scored. Together they are enough to reproduce the run.}
